\section{盐、寺院与流动的人}
\label{sec:late-tang-money-religion-mobility}

安史以后，盐铁转运和两税法不是附在战争故事后的制度名词。江淮粮、盐利和专使网络使朝廷能继续支付军队、维持两京，也使收入节点成为宦官、朝臣与军镇争夺的对象。制度志把这些变化整理为改革谱系，敦煌等地文书则让人看见某些地区的手实、计帐和契约；二者都不能直接给出全国收入或普通家庭负担的总账。%
\footnote{刘晏转运、两税法与盐池冲突见 LT-758-007、LT-780-020、LT-885-067，依据《旧唐书》食货志、刘晏传、杨炎传、王重荣传及《通典》。盐价、税率和留截必须按区域和时间检验。}

819年佛骨入宫引起韩愈上表，845年废寺、还俗与没收资产，则显示宗教并非财政与国家权力之外的领域。迎奉、施舍、寺院土地和铜料，都把皇室、僧俗、地方施主与市场关系连在一起。两种事件不能被简化为“崇佛”与“反佛”的性格对立：前者有礼仪和政治声望的语境，后者则发生在财政、铸钱、户籍和军费均紧张的条件下。%
\footnote{佛骨与会昌废佛见 LT-819-037、LT-845-049，依据《旧唐书》宪宗、武宗纪、《资治通鉴》及佛教传记材料。寺院数、还俗人数和资产收取的数字是行政或宗教论辩的一部分，不能不加辨析地全国化。}

回鹘部众迁徙、河西归义军的形成、安南的失守和收复、黄巢军的南北移动，使人看到晚唐并不存在一条单向的“人口流亡”路线。军队、商人、僧侣、逃户和地方精英各以不同方式使用道路、港口、绿洲与城墙；史书多在其威胁京师时才留下强烈叙事。把敦煌文书、港市遗存与中原正史并读，能让我们提出迁徙和供给的问题，不能替这些材料没有保存的人群编造统一经验。%
\footnote{回鹘迁徙、归义军、安南与黄巢广州行军见 LT-840-045、LT-848-051、LT-851-052、LT-863-054、LT-866-055、LT-879-061。可参《旧唐书》回鹘、张议潮、高骈、黄巢诸传、敦煌文书及港口考古；各类材料的空间覆盖极不均衡。}
